% --- Archon physics formalization source begin ---
% archon:physics
% archon:covers PhyXMiniProblems/problem_phyx_mini_0881.lean
% archon:source-report reports/phyx_mini/problem_phyx_mini_0881.source.json

\chapter{Physics problem phyx\_mini\_0881}
\label{ch:PhyXMiniProblems_problem_phyx_mini_0881}

\paragraph{Problem source.}
\#\# Physical scenario

The switch in figure has been open for a long time. It is closed at \textbackslash{}( t = 0 \textbackslash{}, \textbackslash{}text\{s\} \textbackslash{}).

\#\# Question

What is the current through the battery after the switch has been closed a long time?

\#\# Answer choices

- A: A: \textbackslash{}( 0.0 \textbackslash{}

- B: B: \textbackslash{}( 4.0 \textbackslash{}

- C: C: \textbackslash{}( 1.0 \textbackslash{}

- D: D: \textbackslash{}( 0.5 \textbackslash{}

\#\# Figure caption (auxiliary; use the image as primary evidence)

The image depicts an electrical circuit diagram consisting of the following components:

1. **Voltage Source**: A 10 V battery is present, marked with positive (+) and negative (-) terminals.

2. **Switch**: A switch is shown in the open position, indicating that the circuit is initially incomplete.

3. **Components in Parallel**:
   - Two resistors are present in parallel configuration.
   - Each resistor is labeled with a resistance of 20 Ω.

4. **Inductor**: There is an inductor with an inductance of 10 mH in series with one of the resistors.

The circuit layout suggests a combination of series and parallel components. The switch is placed in series with the source and the parallel configuration of the resistor and the inductor-resistor pair. The lines represent conductive traces or wires connecting the components.

\#\# Dataset metadata

Electric Circuits; Physical Model Grounding Reasoning, Implicit Condition Reasoning

\paragraph{Recorded answer/context.}
C: C: \textbackslash{}( 1.0 \textbackslash{}

\paragraph{Figure/image path.}
/root/proposal\_for\_physic/science-mango/phyx\_mini\_run/phyx\_data/test\_image/881.png

\paragraph{Formalization target.}
create a compiling Lean file with sorry bodies at `PhyXMiniProblems/problem\_phyx\_mini\_0881.lean`. The Lean declarations must preserve the physical quantities, dimensions or dimensional roles, figure labels, governing-law hypotheses, and final relation expressed by this problem.
Use Mathlib/Physlib names found through LeanExplore where available. If a physics API is missing, introduce faithful local abstractions rather than scalar placeholder aliases.

\begin{theorem}[Physics formalization target]
\label{thm:physics:phyx_mini_0881:target}
\lean{PhyXMiniProblems.ProblemPhyXMini0881.problem_phyx_mini_0881}
\uses{def:physics:phyx-mini-0881:phyxminiproblems-problemphyxmini0881-currentinamperes, def:physics:phyx-mini-0881:phyxminiproblems-problemphyxmini0881-switchedparallelrlcircuit, def:physics:phyx-mini-0881:phyxminiproblems-problemphyxmini0881-matchessuppliedparallelrlfigure, def:physics:phyx-mini-0881:phyxminiproblems-problemphyxmini0881-matchesswitchingprotocol, def:physics:phyx-mini-0881:phyxminiproblems-problemphyxmini0881-hasphysicalcircuitparameters, def:physics:phyx-mini-0881:phyxminiproblems-problemphyxmini0881-satisfieslongtimeidealdccircuitlaws, def:physics:phyx-mini-0881:phyxminiproblems-problemphyxmini0881-isuniquematchinganswer}
The assigned autoformalize agent should translate this physics problem into Lean declarations in the covered file, with theorem and lemma proof bodies written as `by sorry`.
\end{theorem}
\begin{proof}
This is an autoformalization task, not a proof task. Produce statements that can later be proved by the physics prover without weakening the physical model.
\end{proof}

% --- Archon declaration topology begin ---
\section{Declaration topology}
\begin{definition}[energy Dimension]
  \label{def:physics:phyx-mini-0881:phyxminiproblems-problemphyxmini0881-energydimension}
  \lean{PhyXMiniProblems.ProblemPhyXMini0881.energyDimension}
  Energy has dimension M L² T⁻²
\end{definition}
\begin{proof}
  This is the declared model definition of energy Dimension.
\end{proof}
\begin{definition}[electric Current Dimension]
  \label{def:physics:phyx-mini-0881:phyxminiproblems-problemphyxmini0881-electriccurrentdimension}
  \lean{PhyXMiniProblems.ProblemPhyXMini0881.electricCurrentDimension}
  Electric current has dimension charge per time
\end{definition}
\begin{proof}
  This is the declared model definition of electric Current Dimension.
\end{proof}
\begin{definition}[electric Potential Dimension]
  \label{def:physics:phyx-mini-0881:phyxminiproblems-problemphyxmini0881-electricpotentialdimension}
  \lean{PhyXMiniProblems.ProblemPhyXMini0881.electricPotentialDimension}
  \uses{def:physics:phyx-mini-0881:phyxminiproblems-problemphyxmini0881-energydimension}
  Electric potential difference has dimension energy per charge
\end{definition}
\begin{proof}
  This is the declared model definition of electric Potential Dimension.
\end{proof}
\begin{definition}[electrical Resistance Dimension]
  \label{def:physics:phyx-mini-0881:phyxminiproblems-problemphyxmini0881-electricalresistancedimension}
  \lean{PhyXMiniProblems.ProblemPhyXMini0881.electricalResistanceDimension}
  \uses{def:physics:phyx-mini-0881:phyxminiproblems-problemphyxmini0881-electriccurrentdimension, def:physics:phyx-mini-0881:phyxminiproblems-problemphyxmini0881-electricpotentialdimension}
  Electrical resistance has dimension potential difference per current
\end{definition}
\begin{proof}
  This is the declared model definition of electrical Resistance Dimension.
\end{proof}
\begin{definition}[inductance Dimension]
  \label{def:physics:phyx-mini-0881:phyxminiproblems-problemphyxmini0881-inductancedimension}
  \lean{PhyXMiniProblems.ProblemPhyXMini0881.inductanceDimension}
  \uses{def:physics:phyx-mini-0881:phyxminiproblems-problemphyxmini0881-electricalresistancedimension}
  Inductance has dimension resistance times time
\end{definition}
\begin{proof}
  This is the declared model definition of inductance Dimension.
\end{proof}
\begin{definition}[Voltage Quantity]
  \label{def:physics:phyx-mini-0881:phyxminiproblems-problemphyxmini0881-voltagequantity}
  \lean{PhyXMiniProblems.ProblemPhyXMini0881.VoltageQuantity}
  \uses{def:physics:phyx-mini-0881:phyxminiproblems-problemphyxmini0881-electricpotentialdimension}
  A nonnegative, unit-independent voltage magnitude
\end{definition}
\begin{proof}
  This is the declared model definition of Voltage Quantity.
\end{proof}
\begin{definition}[Resistance Quantity]
  \label{def:physics:phyx-mini-0881:phyxminiproblems-problemphyxmini0881-resistancequantity}
  \lean{PhyXMiniProblems.ProblemPhyXMini0881.ResistanceQuantity}
  \uses{def:physics:phyx-mini-0881:phyxminiproblems-problemphyxmini0881-electricalresistancedimension}
  A nonnegative, unit-independent electrical resistance
\end{definition}
\begin{proof}
  This is the declared model definition of Resistance Quantity.
\end{proof}
\begin{definition}[Inductance Quantity]
  \label{def:physics:phyx-mini-0881:phyxminiproblems-problemphyxmini0881-inductancequantity}
  \lean{PhyXMiniProblems.ProblemPhyXMini0881.InductanceQuantity}
  \uses{def:physics:phyx-mini-0881:phyxminiproblems-problemphyxmini0881-inductancedimension}
  A nonnegative, unit-independent inductance
\end{definition}
\begin{proof}
  This is the declared model definition of Inductance Quantity.
\end{proof}
\begin{definition}[Electric Current Quantity]
  \label{def:physics:phyx-mini-0881:phyxminiproblems-problemphyxmini0881-electriccurrentquantity}
  \lean{PhyXMiniProblems.ProblemPhyXMini0881.ElectricCurrentQuantity}
  \uses{def:physics:phyx-mini-0881:phyxminiproblems-problemphyxmini0881-electriccurrentdimension}
  A nonnegative, unit-independent current magnitude in a fixed orientation
\end{definition}
\begin{proof}
  This is the declared model definition of Electric Current Quantity.
\end{proof}
\begin{definition}[coherent SIReadout]
  \label{def:physics:phyx-mini-0881:phyxminiproblems-problemphyxmini0881-coherentsireadout}
  \lean{PhyXMiniProblems.ProblemPhyXMini0881.coherentSIReadout}
  Coherent-SI scalar readout of a nonnegative dimensionful quantity
\end{definition}
\begin{proof}
  This is the declared model definition of coherent SIReadout.
\end{proof}
\begin{definition}[voltage In Volts]
  \label{def:physics:phyx-mini-0881:phyxminiproblems-problemphyxmini0881-voltageinvolts}
  \lean{PhyXMiniProblems.ProblemPhyXMini0881.voltageInVolts}
  \uses{def:physics:phyx-mini-0881:phyxminiproblems-problemphyxmini0881-voltagequantity, def:physics:phyx-mini-0881:phyxminiproblems-problemphyxmini0881-coherentsireadout}
  Read a voltage magnitude in volts
\end{definition}
\begin{proof}
  This is the declared model definition of voltage In Volts.
\end{proof}
\begin{definition}[resistance In Ohms]
  \label{def:physics:phyx-mini-0881:phyxminiproblems-problemphyxmini0881-resistanceinohms}
  \lean{PhyXMiniProblems.ProblemPhyXMini0881.resistanceInOhms}
  \uses{def:physics:phyx-mini-0881:phyxminiproblems-problemphyxmini0881-resistancequantity, def:physics:phyx-mini-0881:phyxminiproblems-problemphyxmini0881-coherentsireadout}
  Read an electrical resistance in ohms
\end{definition}
\begin{proof}
  This is the declared model definition of resistance In Ohms.
\end{proof}
\begin{definition}[inductance In Henries]
  \label{def:physics:phyx-mini-0881:phyxminiproblems-problemphyxmini0881-inductanceinhenries}
  \lean{PhyXMiniProblems.ProblemPhyXMini0881.inductanceInHenries}
  \uses{def:physics:phyx-mini-0881:phyxminiproblems-problemphyxmini0881-inductancequantity, def:physics:phyx-mini-0881:phyxminiproblems-problemphyxmini0881-coherentsireadout}
  Read an inductance in henries
\end{definition}
\begin{proof}
  This is the declared model definition of inductance In Henries.
\end{proof}
\begin{definition}[inductance In Millihenries]
  \label{def:physics:phyx-mini-0881:phyxminiproblems-problemphyxmini0881-inductanceinmillihenries}
  \lean{PhyXMiniProblems.ProblemPhyXMini0881.inductanceInMillihenries}
  \uses{def:physics:phyx-mini-0881:phyxminiproblems-problemphyxmini0881-inductancequantity, def:physics:phyx-mini-0881:phyxminiproblems-problemphyxmini0881-inductanceinhenries}
  Read an inductance in millihenries
\end{definition}
\begin{proof}
  This is the declared model definition of inductance In Millihenries.
\end{proof}
\begin{definition}[current In Amperes]
  \label{def:physics:phyx-mini-0881:phyxminiproblems-problemphyxmini0881-currentinamperes}
  \lean{PhyXMiniProblems.ProblemPhyXMini0881.currentInAmperes}
  \uses{def:physics:phyx-mini-0881:phyxminiproblems-problemphyxmini0881-electriccurrentquantity, def:physics:phyx-mini-0881:phyxminiproblems-problemphyxmini0881-coherentsireadout}
  Read a current magnitude in amperes
\end{definition}
\begin{proof}
  This is the declared model definition of current In Amperes.
\end{proof}
\begin{definition}[Resistor Position]
  \label{def:physics:phyx-mini-0881:phyxminiproblems-problemphyxmini0881-resistorposition}
  \lean{PhyXMiniProblems.ProblemPhyXMini0881.ResistorPosition}
  The two resistor positions distinguished by the supplied circuit diagram
\end{definition}
\begin{proof}
  This is the declared model definition of Resistor Position.
\end{proof}
\begin{definition}[Circuit Component]
  \label{def:physics:phyx-mini-0881:phyxminiproblems-problemphyxmini0881-circuitcomponent}
  \lean{PhyXMiniProblems.ProblemPhyXMini0881.CircuitComponent}
  Individually identifiable components visible in image 881.png
\end{definition}
\begin{proof}
  This is the declared model definition of Circuit Component.
\end{proof}
\begin{definition}[Circuit Topology]
  \label{def:physics:phyx-mini-0881:phyxminiproblems-problemphyxmini0881-circuittopology}
  \lean{PhyXMiniProblems.ProblemPhyXMini0881.CircuitTopology}
  The branch connectivity shown in the primary raster
\end{definition}
\begin{proof}
  This is the declared model definition of Circuit Topology.
\end{proof}
\begin{definition}[Switch State]
  \label{def:physics:phyx-mini-0881:phyxminiproblems-problemphyxmini0881-switchstate}
  \lean{PhyXMiniProblems.ProblemPhyXMini0881.SwitchState}
  State of the switch as a function of the time coordinate
\end{definition}
\begin{proof}
  This is the declared model definition of Switch State.
\end{proof}
\begin{definition}[Parallel RLFigure]
  \label{def:physics:phyx-mini-0881:phyxminiproblems-problemphyxmini0881-parallelrlfigure}
  \lean{PhyXMiniProblems.ProblemPhyXMini0881.ParallelRLFigure}
  \uses{def:physics:phyx-mini-0881:phyxminiproblems-problemphyxmini0881-resistorposition, def:physics:phyx-mini-0881:phyxminiproblems-problemphyxmini0881-circuitcomponent, def:physics:phyx-mini-0881:phyxminiproblems-problemphyxmini0881-circuittopology}
  Literal presentation data from the supplied figure. Optional scalar fields record whether a numerical label is printed; they do not define any physical quantity in the circuit setup
\end{definition}
\begin{proof}
  This is the declared model definition of Parallel RLFigure.
\end{proof}
\begin{definition}[Switched Parallel RLCircuit]
  \label{def:physics:phyx-mini-0881:phyxminiproblems-problemphyxmini0881-switchedparallelrlcircuit}
  \lean{PhyXMiniProblems.ProblemPhyXMini0881.SwitchedParallelRLCircuit}
  \uses{def:physics:phyx-mini-0881:phyxminiproblems-problemphyxmini0881-voltagequantity, def:physics:phyx-mini-0881:phyxminiproblems-problemphyxmini0881-resistancequantity, def:physics:phyx-mini-0881:phyxminiproblems-problemphyxmini0881-inductancequantity, def:physics:phyx-mini-0881:phyxminiproblems-problemphyxmini0881-electriccurrentquantity, def:physics:phyx-mini-0881:phyxminiproblems-problemphyxmini0881-switchstate, def:physics:phyx-mini-0881:phyxminiproblems-problemphyxmini0881-parallelrlfigure}
  Independent physical quantities for the switched circuit. The functions of real time use seconds as their coordinate. Long-time currents and voltage drops are fields rather than definitions from the displayed answer
\end{definition}
\begin{proof}
  This is the declared model definition of Switched Parallel RLCircuit.
\end{proof}
\begin{definition}[Matches Supplied Parallel RLFigure]
  \label{def:physics:phyx-mini-0881:phyxminiproblems-problemphyxmini0881-matchessuppliedparallelrlfigure}
  \lean{PhyXMiniProblems.ProblemPhyXMini0881.MatchesSuppliedParallelRLFigure}
  \uses{def:physics:phyx-mini-0881:phyxminiproblems-problemphyxmini0881-voltageinvolts, def:physics:phyx-mini-0881:phyxminiproblems-problemphyxmini0881-resistanceinohms, def:physics:phyx-mini-0881:phyxminiproblems-problemphyxmini0881-inductanceinmillihenries, def:physics:phyx-mini-0881:phyxminiproblems-problemphyxmini0881-switchedparallelrlcircuit}
  Primary-raster evidence and its calibration to independent dimensionful quantities. This records the 10 V, two 20 Ω, and 10 mH labels together with the depicted branch connectivity and polarity marks
\end{definition}
\begin{proof}
  This is the declared model definition of Matches Supplied Parallel RLFigure.
\end{proof}
\begin{definition}[Matches Switching Protocol]
  \label{def:physics:phyx-mini-0881:phyxminiproblems-problemphyxmini0881-matchesswitchingprotocol}
  \lean{PhyXMiniProblems.ProblemPhyXMini0881.MatchesSwitchingProtocol}
  \uses{def:physics:phyx-mini-0881:phyxminiproblems-problemphyxmini0881-switchedparallelrlcircuit}
  The switch is idealized as having been open throughout negative times, being closed at t = 0 s, and remaining closed thereafter
\end{definition}
\begin{proof}
  This is the declared model definition of Matches Switching Protocol.
\end{proof}
\begin{definition}[Has Physical Circuit Parameters]
  \label{def:physics:phyx-mini-0881:phyxminiproblems-problemphyxmini0881-hasphysicalcircuitparameters}
  \lean{PhyXMiniProblems.ProblemPhyXMini0881.HasPhysicalCircuitParameters}
  \uses{def:physics:phyx-mini-0881:phyxminiproblems-problemphyxmini0881-voltageinvolts, def:physics:phyx-mini-0881:phyxminiproblems-problemphyxmini0881-resistanceinohms, def:physics:phyx-mini-0881:phyxminiproblems-problemphyxmini0881-inductanceinhenries, def:physics:phyx-mini-0881:phyxminiproblems-problemphyxmini0881-switchedparallelrlcircuit}
  Positivity conditions for the three passive component parameters
\end{definition}
\begin{proof}
  This is the declared model definition of Has Physical Circuit Parameters.
\end{proof}
\begin{definition}[Satisfies Long Time Ideal DCCircuit Laws]
  \label{def:physics:phyx-mini-0881:phyxminiproblems-problemphyxmini0881-satisfieslongtimeidealdccircuitlaws}
  \lean{PhyXMiniProblems.ProblemPhyXMini0881.SatisfiesLongTimeIdealDCCircuitLaws}
  \uses{def:physics:phyx-mini-0881:phyxminiproblems-problemphyxmini0881-voltageinvolts, def:physics:phyx-mini-0881:phyxminiproblems-problemphyxmini0881-resistanceinohms, def:physics:phyx-mini-0881:phyxminiproblems-problemphyxmini0881-currentinamperes, def:physics:phyx-mini-0881:phyxminiproblems-problemphyxmini0881-switchedparallelrlcircuit}
  The ideal lumped-circuit laws used at long times: an open switch disconnects the battery; the time-dependent battery current tends to the named long-time current; parallel branches share the battery voltage, and Kirchhoff's voltage law splits the series-branch voltage between its resistor and inductor; both resistors obey Ohm's law; an ideal inductor in the long-time DC regime has zero voltage drop; and Kirchhoff's current law adds the two branch currents at the battery.
\end{definition}
\begin{proof}
  This is the declared model definition of Satisfies Long Time Ideal DCCircuit Laws.
\end{proof}
\begin{lemma}[long Time branch Currents eq one Half Ampere]
  \label{lem:physics:phyx-mini-0881:phyxminiproblems-problemphyxmini0881-longtime-branchcurrents-eq-onehalfampere}
  \lean{PhyXMiniProblems.ProblemPhyXMini0881.longTime_branchCurrents_eq_oneHalfAmpere}
  \uses{def:physics:phyx-mini-0881:phyxminiproblems-problemphyxmini0881-currentinamperes, def:physics:phyx-mini-0881:phyxminiproblems-problemphyxmini0881-switchedparallelrlcircuit, def:physics:phyx-mini-0881:phyxminiproblems-problemphyxmini0881-matchessuppliedparallelrlfigure, def:physics:phyx-mini-0881:phyxminiproblems-problemphyxmini0881-satisfieslongtimeidealdccircuitlaws}
  Each 20 Ω branch carries 0.5 A in the long-time DC regime: the ideal inductor has become a zero-voltage series element, so each resistor has the full 10 V source voltage across it
\end{lemma}
\begin{proof}
  The conclusion follows from the governing relations and figure readouts named in the statement of long Time branch Currents eq one Half Ampere.
\end{proof}
\begin{definition}[Answer Choice]
  \label{def:physics:phyx-mini-0881:phyxminiproblems-problemphyxmini0881-answerchoice}
  \lean{PhyXMiniProblems.ProblemPhyXMini0881.AnswerChoice}
  Labels of the four current choices printed with the source problem
\end{definition}
\begin{proof}
  This is the declared model definition of Answer Choice.
\end{proof}
\begin{definition}[displayed Current In Amperes]
  \label{def:physics:phyx-mini-0881:phyxminiproblems-problemphyxmini0881-displayedcurrentinamperes}
  \lean{PhyXMiniProblems.ProblemPhyXMini0881.displayedCurrentInAmperes}
  \uses{def:physics:phyx-mini-0881:phyxminiproblems-problemphyxmini0881-answerchoice}
  Current in amperes printed beside each answer label
\end{definition}
\begin{proof}
  This is the declared model definition of displayed Current In Amperes.
\end{proof}
\begin{definition}[recorded Dataset Answer]
  \label{def:physics:phyx-mini-0881:phyxminiproblems-problemphyxmini0881-recordeddatasetanswer}
  \lean{PhyXMiniProblems.ProblemPhyXMini0881.recordedDatasetAnswer}
  \uses{def:physics:phyx-mini-0881:phyxminiproblems-problemphyxmini0881-answerchoice}
  The source dataset's recorded answer label, retained only as metadata
\end{definition}
\begin{proof}
  This is the declared model definition of recorded Dataset Answer.
\end{proof}
\begin{definition}[Matches Long Time Battery Current]
  \label{def:physics:phyx-mini-0881:phyxminiproblems-problemphyxmini0881-matcheslongtimebatterycurrent}
  \lean{PhyXMiniProblems.ProblemPhyXMini0881.MatchesLongTimeBatteryCurrent}
  \uses{def:physics:phyx-mini-0881:phyxminiproblems-problemphyxmini0881-currentinamperes, def:physics:phyx-mini-0881:phyxminiproblems-problemphyxmini0881-switchedparallelrlcircuit, def:physics:phyx-mini-0881:phyxminiproblems-problemphyxmini0881-answerchoice, def:physics:phyx-mini-0881:phyxminiproblems-problemphyxmini0881-displayedcurrentinamperes}
  A displayed choice equals the modeled long-time battery current
\end{definition}
\begin{proof}
  This is the declared model definition of Matches Long Time Battery Current.
\end{proof}
\begin{definition}[Is Unique Matching Answer]
  \label{def:physics:phyx-mini-0881:phyxminiproblems-problemphyxmini0881-isuniquematchinganswer}
  \lean{PhyXMiniProblems.ProblemPhyXMini0881.IsUniqueMatchingAnswer}
  \uses{def:physics:phyx-mini-0881:phyxminiproblems-problemphyxmini0881-switchedparallelrlcircuit, def:physics:phyx-mini-0881:phyxminiproblems-problemphyxmini0881-answerchoice, def:physics:phyx-mini-0881:phyxminiproblems-problemphyxmini0881-matcheslongtimebatterycurrent}
  A displayed choice is the unique exact match for the modeled current
\end{definition}
\begin{proof}
  This is the declared model definition of Is Unique Matching Answer.
\end{proof}
% --- Archon declaration topology end ---
% --- Archon physics formalization source end ---
